\documentclass[11pt]{article}

\usepackage[margin=1in]{geometry}
\usepackage{amsmath,amssymb,amsthm}
\usepackage{graphicx}

\newcommand{\url}[1]{\texttt{#1}}
\newcommand{\toprule}{\hline\hline}
\newcommand{\midrule}{\hline}
\newcommand{\bottomrule}{\hline\hline}

\newtheorem{lemma}{Lemma}
\newtheorem{proposition}{Proposition}

\title{An Approximate Capacity Bound for Sparse Combinatorial Key-Value Memory:\\
Full Derivation and Validation}

\author{Supplementary note to ``Combinatorial Coding, Key-Value Recall, and Emergent\\
Constraint Satisfaction in a Hippocampal-Inspired Spiking Network''}

\date{\today}

\begin{document}
\maketitle

\begin{abstract}
This note gives the complete derivation, with every intermediate step made explicit, of the
approximate capacity bound used in the main paper (Section 2.3): an expression for the largest
number of items $M$ a sparse $K$-of-$H$ combinatorial key-value memory can store while
remaining reliably retrievable under a given level of query corruption. The bound is
structurally the same as modern Hopfield networks' proven exponential capacity
\cite{ramsauer2020}, but is derived here for a discrete, sparse code via elementary
combinatorics and extreme value theory rather than for dense continuous patterns via an
energy-function argument. We include a full, self-contained proof of the one non-elementary
step (the expected maximum of many correlated-in-effect-but-treated-as-independent Gaussian
similarities), validate the final bound against simulation, and discuss its relationship to
the wider family of capacity theorems it belongs to.
\end{abstract}

\section{Setup}

We store $M$ items. Item $m$'s key is a binary vector $h_m \in \{0,1\}^H$ with exactly $K$
active (value-$1$) coordinates -- a $K$-of-$H$ sparse code. Throughout this note we model the
$M$ stored keys as independent, uniformly random $K$-subsets of $[H] = \{1,\dots,H\}$: an
idealization relative to the true architecture, where keys are $W_{SH}$-projections of
mutually correlated combinatorial structure (solved Sudoku grids), addressed explicitly in
Section~\ref{sec:limitations}.

A query is formed from the true key $h_{\text{true}}$ by flipping a fraction $\eta \in [0,1]$
of all $H$ coordinates, chosen uniformly at random without replacement -- the noise model used
throughout the main paper. Retrieval compares the query against all $M$ stored keys via
similarity (shared active coordinates) and should return the index of $h_{\text{true}}$.
Retrieval \emph{fails} when some other stored key's similarity to the query exceeds
$h_{\text{true}}$'s.

We ask: for fixed $H$, $K$, and corruption level $\eta$, how large can $M$ be before this kind
of failure becomes likely?

\section{Step 1: Similarity between two independent random keys}
\label{sec:step1}

\begin{lemma}[Hypergeometric overlap]
\label{lem:hypergeom}
Let $A, B \subset [H]$ be independent uniformly random subsets with $|A| = a$, $|B| = b$. Then
$|A \cap B|$ is a hypergeometric random variable with
\[
\mathbb{E}[|A \cap B|] = \frac{ab}{H}, \qquad
\mathrm{Var}(|A \cap B|) = a\cdot\frac{b}{H}\Big(1-\frac{b}{H}\Big)\cdot\frac{H-a}{H-1}.
\]
\end{lemma}

\begin{proof}[Proof sketch]
Fix $A$ and treat $B$'s $b$ elements as drawn sequentially without replacement from $[H]$.
Let $X_i = \mathbb{1}[\text{the $i$-th draw lands in } A]$, so $|A\cap B| = \sum_{i=1}^b X_i$.
Each $X_i$ is Bernoulli$(a/H)$ by symmetry, giving the mean by linearity of expectation. For
the variance, $\mathrm{Cov}(X_i,X_j) < 0$ for $i\ne j$ because sampling without replacement
makes draws negatively associated (having drawn a marked element makes drawing another
marked element on a later draw slightly less likely); summing the $b$ variances and
$\binom{b}{2}$ covariances gives the stated finite-population-corrected formula. (Full
derivation: Feller, \textit{An Introduction to Probability Theory}, Vol.\ 1, Ch.\ II.)
\end{proof}

Applying Lemma~\ref{lem:hypergeom} with $a=b=K$ gives the similarity between two
\emph{uncorrupted} stored keys: mean $K^2/H$, the number quoted throughout the main paper (e.g.\
$K^2/H = 22.5$ for the default $H=1000$, $K=150$).

\section{Step 2: Statistics of a corrupted query}
\label{sec:step2}

Flipping a uniformly random fraction $\eta$ of all $H$ coordinates is itself an instance of
Lemma~\ref{lem:hypergeom}: the $\eta H$ flipped positions are a random subset of $[H]$, and we
ask how many of $h_{\text{true}}$'s $K$ active positions they contain. By the lemma (with
$a=K$, $b=\eta H$), the expected number of active positions flipped \emph{off} is $\eta K$, so
\begin{equation}
S_{\text{true}}(\eta) \;=\; \mathbb{E}\big[\,\text{similarity of query to } h_{\text{true}}\,\big] \;=\; K(1-\eta),
\label{eq:strue}
\end{equation}
exactly (no further approximation needed: this is a direct expectation, not a bound). The
expected number of originally-\emph{inactive} positions flipped \emph{on} is, symmetrically,
$\eta(H-K)$, so the corrupted query's own total active-coordinate count is, in expectation,
\begin{equation}
K_q(\eta) \;=\; K(1-\eta) + \eta(H-K) \;=\; K(1-2\eta) + \eta H.
\label{eq:kq}
\end{equation}
The query's similarity to any \emph{other}, independent stored key is now Lemma~\ref{lem:hypergeom}
applied with $a = K_q(\eta)$, $b=K$:
\begin{equation}
\mu_{\text{other}}(\eta) = \frac{K_q(\eta)\,K}{H}, \qquad
\sigma_{\text{other}}^2(\eta) = K_q(\eta)\,\frac{K}{H}\Big(1-\frac{K}{H}\Big)\frac{H-K_q(\eta)}{H-1}.
\label{eq:muother}
\end{equation}

\section{Step 3: Gaussian approximation}
\label{sec:step3}

The hypergeometric distribution of $S_{\text{other}}$ (similarity to one specific other key) is
exact but discrete and unwieldy for the extreme-value step that follows. We approximate
\[
S_{\text{other}} \;\approx\; \mathcal{N}\big(\mu_{\text{other}}(\eta),\, \sigma_{\text{other}}^2(\eta)\big),
\]
justified for $H \gg 1$ by the central limit theorem for sampling without replacement from a
finite population \cite{erdosrenyi1959,hajek1960}: sums arising from simple random sampling
without replacement are asymptotically normal as the population size grows, with exactly the
finite-population-corrected variance already appearing in Lemma~\ref{lem:hypergeom}. This is
the least exact step in the derivation, and its error is largest precisely when $\sigma_{\text{other}}$
is small relative to the mean -- i.e.\ at low noise, where the discreteness of the underlying
hypergeometric distribution matters most. Section~\ref{sec:validation} quantifies this error
directly rather than leaving it as a caveat.

\section{Step 4: The maximum of many competitors}
\label{sec:step4}

Retrieval does not fail when the \emph{typical} competing key looks similar to the query --
it fails when the \emph{single most similar} of the $M-1$ competitors exceeds
$S_{\text{true}}(\eta)$. This is an extreme-value criterion, and it is the step that most
distinguishes this derivation from a naive signal-to-noise argument.

\begin{lemma}[Expected maximum of i.i.d.\ Gaussians]
\label{lem:extreme}
Let $Z_1,\dots,Z_n$ be i.i.d.\ $\mathcal{N}(0,1)$. Then
\[
\mathbb{E}\Big[\max_{1\le i\le n} Z_i\Big] \;\le\; \sqrt{2\ln n} \;+\; o(1) \quad \text{as } n\to\infty,
\]
and this bound is asymptotically tight: $\mathbb{E}[\max_i Z_i] = \sqrt{2\ln n}\,(1+o(1))$.
\end{lemma}

\begin{proof}[Proof of the upper bound]
For any $t \ge 0$, the standard Gaussian tail (Chernoff) bound gives $\Pr(Z_1 > t) \le
e^{-t^2/2}$: with $s=t$, $\Pr(Z_1>t) = \Pr(e^{sZ_1} > e^{st}) \le \mathbb{E}[e^{sZ_1}]e^{-st} =
e^{s^2/2 - st}$, minimized at $s=t$. A union bound over the $n$ variables gives $\Pr(\max_i Z_i
> t) \le \min(1,\, n e^{-t^2/2})$. Writing the expectation of a nonnegative random variable as
an integral of its tail,
\[
\mathbb{E}\Big[\max_i Z_i\Big] = \int_0^\infty \Pr\Big(\max_i Z_i > t\Big)\,dt
\;\le\; \int_0^{t^*} 1 \,dt + \int_{t^*}^\infty n e^{-t^2/2}\,dt,
\]
where $t^* = \sqrt{2\ln n}$ is chosen so the two bounds on the integrand meet ($n e^{-(t^*)^2/2}
= 1$). The first integral is exactly $t^* = \sqrt{2\ln n}$. For the second, the standard
Gaussian tail estimate $\int_a^\infty e^{-t^2/2}\,dt \le \tfrac{1}{a}e^{-a^2/2}$ (for $a>0$,
obtained by writing $e^{-t^2/2} \le \tfrac{t}{a}e^{-t^2/2}$ on $t\ge a$ and integrating exactly)
gives
\[
\int_{t^*}^\infty n e^{-t^2/2}\,dt \;\le\; \frac{n}{t^*}e^{-(t^*)^2/2} = \frac{1}{\sqrt{2\ln n}} \;\xrightarrow{n\to\infty}\; 0.
\]
Summing the two pieces gives $\mathbb{E}[\max_i Z_i] \le \sqrt{2\ln n} + o(1)$. The matching
lower bound requires a second-moment (Paley--Zygmund-type) argument and is omitted here; see
Boucheron, Lugosi \& Massart, \textit{Concentration Inequalities} \cite{blm2013}, Ch.\ 2, for
the complete two-sided result.
\end{proof}

Rescaling Lemma~\ref{lem:extreme} to $\mathcal{N}(\mu_{\text{other}},\sigma_{\text{other}}^2)$
variables and setting $n = M-1$ (the number of competing stored keys) gives the working
approximation used in the main paper:
\begin{equation}
\mathbb{E}\Big[\max_{i=1}^{M-1} S_{\text{other},i}\Big] \;\approx\; \mu_{\text{other}}(\eta) + \sigma_{\text{other}}(\eta)\sqrt{2\ln(M-1)}.
\label{eq:expmax}
\end{equation}

\section{Step 5: Solving for capacity}
\label{sec:step5}

Retrieval is (approximately, at the level of expectations rather than a proper tail bound on
failure probability) reliable when $S_{\text{true}}(\eta)$ exceeds the expected maximum
competitor, Equation~\ref{eq:expmax}. Define the normalized separation margin
\begin{equation}
z(\eta) \;=\; \frac{S_{\text{true}}(\eta) - \mu_{\text{other}}(\eta)}{\sigma_{\text{other}}(\eta)}.
\label{eq:margin}
\end{equation}
Setting $S_{\text{true}}(\eta) = \mu_{\text{other}}(\eta) + \sigma_{\text{other}}(\eta)\sqrt{2\ln(M-1)}$
and solving for $M$:
\begin{equation}
z(\eta) = \sqrt{2\ln(M-1)} \quad\Longleftrightarrow\quad M-1 = \exp\!\left(\frac{z(\eta)^2}{2}\right)
\quad\Longleftrightarrow\quad
\boxed{M_{\max}(\eta) \;\approx\; 1 + \exp\!\left(\frac{z(\eta)^2}{2}\right).}
\label{eq:mmax}
\end{equation}

\begin{proposition}[Structural form]
Equation~\ref{eq:mmax} states that capacity is exponential in the \emph{square} of a
normalized signal-to-noise-like quantity. This is the same structural form as modern Hopfield
networks' proven capacity $C \cong 2^{d/2}$ \cite{ramsauer2020}, though that bound is derived
via a Lyapunov energy-function argument (the Concave--Convex Procedure applied to a
log-sum-exp energy) over dense continuous patterns, not via hypergeometric statistics over a
sparse discrete code.
\end{proposition}

\section{Validation against simulation}
\label{sec:validation}

Equations \ref{eq:strue}--\ref{eq:mmax} were implemented exactly as derived
(\texttt{capacity\_bound\_derivation.py}, no free parameters fit to data) and compared against
the empirically measured collapse thresholds reported in the main paper.

\begin{center}
\begin{tabular}{lccccc}
\toprule
 & $M{=}50$ & $M{=}1000$ & $M{=}10{,}000$ & $K{=}10$ & $K{=}150$ \\
\midrule
Predicted (this note, $H{=}1000$) & 43.9\% & 41.9\% & 40.6\% & 24.5\% & 41.9\% \\
Observed (simulation) & 45\% & 45\% & 40\% & 25\% & 40--45\% \\
\bottomrule
\end{tabular}
\end{center}

Agreement is within 1--5 percentage points across a 200-fold range of $M$ and an order of
magnitude in $K$, including an almost exact match at $K=10$. A second, independent check: at
$H{=}1000$, $K{=}150$, $M{=}200$, the derivation predicts the ratio $S_{\text{true}}(\eta) /
\mathbb{E}[\max_{i} S_{\text{other},i}]$ at $\eta \in \{0, 0.20, 0.30\}$ as $4.21$, $1.99$,
$1.47$ respectively, against independently measured ratios of $3.39$, $1.92$, $1.45$ (obtained
from a completely separate experiment, the separation-operator ablation, not used to fit any
parameter of this derivation). The two noisy points match to within 4\% and 1.4\%; the
noise-free point is over-predicted by about 24\%, consistent with Section~\ref{sec:step3}'s
observation that the Gaussian approximation is weakest exactly there.

\section{Limitations}
\label{sec:limitations}

Two assumptions are false by construction and are stated as such rather than absorbed silently.

\textbf{Independence.} Real stored keys are $W_{SH}$-projections of genuinely correlated
combinatorial objects (solved Sudoku grids share row/column/box structure), not independent
uniform random subsets. The observed collapse thresholds sit consistently at or slightly above
the predicted ones (Section~\ref{sec:validation}); a plausible explanation is that this real
correlation makes the worst-case competing key marginally less extreme than the independent
model assumes, though this has not been tested directly.

\textbf{The extreme-value step is a first-moment approximation, not a tail bound.} Equation
\ref{eq:expmax} matches the \emph{expected} maximum, not a high-probability bound on it; a
fully rigorous capacity statement would instead bound $\Pr(\max_i S_{\text{other},i} >
S_{\text{true}}(\eta))$ directly (e.g.\ via the same union-bound technique used in the proof of
Lemma~\ref{lem:extreme}, without first taking an expectation) and would very likely give a
somewhat smaller $M_{\max}$ at any fixed, small target failure probability. The bound derived
here should be read as characterizing the right \emph{mechanism and order of magnitude} --
exponential-in-margin capacity, hence a sharp rather than gradual collapse, hence a threshold
only weakly dependent on $M$ -- not as a proof.

\section{Relation to the wider literature}

This derivation belongs to a family of results sharing one underlying question: does a correct
answer stand out against the \emph{extreme} value among exponentially or combinatorially many
near-random alternatives? The same skeleton, worked out through different technical routes,
underlies Shannon's channel coding theorem \cite{coverthomas}, the Random Energy Model in
spin-glass theory \cite{mezardmontanari2009}, Gardner's perceptron capacity
\cite{gardner1988}, and the classical Amit--Gutfreund--Sompolinsky analysis of Hopfield network
capacity \cite{ags1987}. Kanerva's original capacity analysis for sparse distributed memory
\cite{kanerva1988} is the closest existing derivation in spirit to this one -- also a sparse
binary code, also analyzed via combinatorial (there, Hamming-ball volume) arguments -- though
via a different geometric route than the hypergeometric-overlap statistics used here.

\begin{thebibliography}{99}

\bibitem{ramsauer2020} H.~Ramsauer et al., ``Hopfield networks is all you need,'' \textit{arXiv:2008.02217} (2020).

\bibitem{erdosrenyi1959} P.~Erd\H{o}s, A.~R\'enyi, ``On a central limit theorem for samples from a finite population,'' \textit{Publ.\ Math.\ Inst.\ Hungarian Acad.\ Sci.} \textbf{4}, 49--61 (1959).

\bibitem{hajek1960} J.~H\'ajek, ``Limiting distributions in simple random sampling from a finite population,'' \textit{Publ.\ Math.\ Inst.\ Hungarian Acad.\ Sci.} \textbf{5}, 361--374 (1960).

\bibitem{blm2013} S.~Boucheron, G.~Lugosi, P.~Massart, \textit{Concentration Inequalities: A Nonasymptotic Theory of Independence}, Oxford University Press (2013).

\bibitem{coverthomas} T.~M.~Cover, J.~A.~Thomas, \textit{Elements of Information Theory}, Wiley (2006).

\bibitem{mezardmontanari2009} M.~M\'ezard, A.~Montanari, \textit{Information, Physics, and Computation}, Oxford University Press (2009).

\bibitem{gardner1988} E.~Gardner, ``The space of interactions in neural network models,'' \textit{J.\ Phys.\ A} \textbf{21}, 257 (1988).

\bibitem{ags1987} D.~J.~Amit, H.~Gutfreund, H.~Sompolinsky, ``Statistical mechanics of neural networks near saturation,'' \textit{Ann.\ Phys.} \textbf{173}, 30--67 (1987).

\bibitem{kanerva1988} P.~Kanerva, \textit{Sparse Distributed Memory}, MIT Press (1988).

\end{thebibliography}

\end{document}
